\title{MLPs are Hebbians: Constructing Efficient Fact-Storing MLPs for Transformers}

\makeatletter
\renewcommand{\thefootnote}{\fnsymbol{footnote}}
\newcommand{\thanksmark}[1]{\textsuperscript{\@fnsymbol{#1}}}
\newcommand{\corrauthormark}{\kern0.35em\textsuperscript{\@fnsymbol{2}}}
\makeatother

\author{
\textbf{Roberto Garcia}\textsuperscript{1}
  \thanks{Equal first authors; order chosen by coin flip}
  \corrauthormark
\quad
\textbf{Jerry Liu}\textsuperscript{1}
  \thanksmark{1}
  \hspace{-0.2cm}
  \corrauthormark
\quad
\textbf{Ronny Junkins}\textsuperscript{2}
  \thanksmark{1}
\\[0.3ex]
\textbf{Sabri Eyuboglu}\textsuperscript{2}
\quad
\textbf{Atri Rudra}\textsuperscript{3}
\quad
\textbf{Chris R{\'e}}\textsuperscript{2}
\\[1.5ex]
\textsuperscript{1}Institute for Computational \& Mathematical Engineering, Stanford University \\
\textsuperscript{2}Department of Computer Science, Stanford University \\
\textsuperscript{3}Department of Computer Science and Engineering, University at Buffalo
}

\newcommand{\hebbianposttitlefootnotes}{%
  \begingroup
  \renewcommand{\thefootnote}{\fnsymbol{footnote}}%
  \footnotetext[2]{Corresponding authors: \href{mailto:robgarct@stanford.edu}{\texttt{robgarct@stanford.edu}}, \href{mailto:jerrywliu@stanford.edu}{\texttt{jerrywliu@stanford.edu}}.}%
  \footnotetext[3]{Code is released at \url{https://github.com/HazyResearch/hebbian-mlps}.}%
  \endgroup
  \setcounter{footnote}{0}%
  \renewcommand{\thefootnote}{\arabic{footnote}}%
}
